\subsection{The actual projective theta-coordinate closure}
\label{subsec:ftb-theta-coordinate-closure}

We begin the fine Siegel family construction at its algebraic
coordinate step.
Let \(o\hookrightarrow\mathbb C\) be a Dedekind domain, let
\(N'>0\), \(n>0\), and put
\[
 \mathcal I=(\mathbb Z/N'^2\mathbb Z)^{2g},\qquad
                              b=\lfloor N'/n\rfloor .
\]
In the eventual application \(N'=n^j\), so \(b=n^{j-1}\).
For a chosen class \(C\) of the source's marked theta objects
of dimension \(g\) and level \(N'^2\), let \(\mathcal V_C\)
be the set of their actual theta-coordinate vectors
\(v:\mathcal I\to\mathbb C\).
These are the vectors in \(\mathrm{HasTh}\): they are
computed from the given marked torsion points, symmetric
theta-group lifts and the specified line section.
Changing the final local scalar multiplying this coordinate
reading by \(c\in\mathbb C^\times\) multiplies the whole
vector by \(c\); all preceding torsion, lift and section
identities remain unchanged.
Thus \(\mathcal V_C\) is stable under nonzero scalar
multiplication.  This is exactly the source's
\(\mathrm{hasTh\_unit\_mul}\) step.

Define the ideal, with no reduction or closure left implicit,
\[
 I_C=\bigcap_{v\in\mathcal V_C}
       \ker\!\left(o[X_a:a\in\mathcal I]\longrightarrow\mathbb C,
                                      \ X_a\longmapsto v_a\right).
 \tag{M1}\label{eq:ftb-theta-vanishing-ideal}
\]
For a ring map \(\beta:o\to R\), write \(\mathrm{InMM}(v)\)
for the conjunction
\[
 \begin{split}
 &F(v)=0\quad(F\in I_C),\\
 &\text{for every prime }q\subset R\text{ and every }a\in\mathcal I,
          \text{ some }c\in\mathcal I
                         \text{ has }v_{a+bc}\notin q.
 \end{split}
 \tag{M2}\label{eq:ftb-theta-inmm}
\]
All polynomial evaluations use \(\beta\).
The second clause is the literal coset nonvanishing condition;
it is stronger than the assertion that \(v\) is not the zero
vector at each point.

For a line bundle \(M\) with sections \(\vartheta_a\) on a
scheme \(Z\), a point \(x:\operatorname{Spec}R\to Z\) has
\emph{section coordinates} \(v\) if there is a generator
\(\tau\) of \(x^*M\) with
\[
                                x^*\vartheta_a=v_a\tau
                                      \quad(a\in\mathcal I).
 \tag{M3}\label{eq:ftb-theta-section-coordinates}
\]
This includes the actual trivialization of the pulled-back
line.  A generator and a basis of an unrelated coordinate
space are not interchangeable in this definition.

\begin{srlemma}[The reduced flat coordinate scheme]
\label{lem:ftb-theta-mm-pre}
The ideal \(I_C\) is radical and homogeneous, and
\(o[X_a]/I_C\) is torsion free over \(o\).
There are a quasiprojective flat reduced \(o\)-scheme
\(\mathcal M_0\), a line \(M\) on it and sections
\(\vartheta_a\), with the following exact property.
For every \(o\)-algebra \(R\), vectors \(v\) satisfying
\(\mathrm{InMM}(v)\) have a unique point
\(x:\operatorname{Spec}R\to\mathcal M_0\) with section
coordinates \(v\), and every such section-coordinate
vector satisfies \(\mathrm{InMM}\).
\end{srlemma}
\begin{proof}
Radicality follows directly from
\eqref{eq:ftb-theta-vanishing-ideal}, because every target
is a field.
For homogeneity, if \(F\in I_C\), then for every actual
coordinate vector \(v\) the polynomial \(F(cv)\) vanishes
for all \(c\ne0\).
Its coefficient of \(c^d\) is the evaluation of the
degree-\(d\) homogeneous part of \(F\).
An infinite set of roots makes this one-variable
polynomial zero, so every homogeneous part lies in
\(I_C\).
Finally, if \(0\ne r\in o\) and \(rF\in I_C\), evaluate
at every \(v\).
The image of \(r\) in \(\mathbb C\) is nonzero, since the
displayed coefficient embedding is injective.
Thus \(F(v)=0\) for every \(v\), proving torsion freeness.

Construct the homogeneous closed subscheme
\(Z\subset\mathbb P_o^{\mathcal I-1}\) defined by \(I_C\).
Its line \(M_Z\) is the restricted \(\mathcal O(1)\);
write \(\vartheta_a\) for its actual coordinate sections.
The standard affine coordinate calculation gives
\[
 D(\vartheta_i)\simeq
 \operatorname{Spec}
  \bigl(o[X_a]/(I_C+(X_i-1))\bigr).
 \tag{M4}\label{eq:ftb-theta-dehom-chart}
\]
For precision, the map sending \(X_a\) to \(X_a/X_i\) in
\((o[X_a]/I_C)[X_i^{-1}]\) has kernel
\(I_C+(X_i-1)\).
One inclusion follows degree by degree from homogeneity.
For the other, evaluate the localized ring at \(X_i=1\):
this is a well-defined retraction because \(X_i\) is
sent to a unit.
Its composite with the ratio map is the quotient map.
This proves the kernel identity and the chart formula,
including the saturation by powers of \(X_i\) which is
already implicit in a projective chart.
Equivalently, a homogeneous form \(F\) vanishes in this
chart exactly when \(X_i^kF\in I_C\) for some \(k\).

The chart ring in \eqref{eq:ftb-theta-dehom-chart} embeds
in the displayed localization.
It is therefore reduced.
It is also torsion free over \(o\).
Indeed, if a nonzero \(r\in o\) kills a localized fraction
\(F/X_i^m\), then some \(X_i^krF\) lies in \(I_C\);
torsion freeness for \(I_C\) gives \(X_i^kF\in I_C\),
so the fraction is zero.
A torsion-free module over a Dedekind domain is flat.

Now take the explicit finite-intersection open
\[
 \mathcal M_0=
      \bigcap_{a\in\mathcal I}
          \ \bigcup_{c\in\mathcal I}D(\vartheta_{a+bc})
                                           \ \subset Z.
 \tag{M5}\label{eq:ftb-theta-coset-open}
\]
The index sets are finite, so this is open.
It is covered by the reduced flat affine charts just
computed; reducedness and flatness are local.
The ambient \(Z\) is projective over the Noetherian
Dedekind domain \(o\); its open \(\mathcal M_0\) is
quasiprojective.

Here are the full point identities.
If \(v\) satisfies \eqref{eq:ftb-theta-inmm}, its coordinates
generate the unit ideal in \(R\): a maximal ideal
containing all of them would violate the coset condition
at \(a=0\).
The ordinary projective coordinate construction therefore
gives a unique point of projective space with
\eqref{eq:ftb-theta-section-coordinates}.
The equations put it in \(Z\).
At each prime of \(R\), the coset condition and
\eqref{eq:ftb-theta-section-coordinates} put its image in
\eqref{eq:ftb-theta-coset-open}.
It factors uniquely through this open immersion.

Conversely a section-coordinate vector on \(\mathcal M_0\)
satisfies all equations, by the homogeneous ideal of \(Z\).
At a prime \(q\), the image of the corresponding point
belongs to every union in
\eqref{eq:ftb-theta-coset-open}.
The line generator in \eqref{eq:ftb-theta-section-coordinates}
is a unit in each stalk; hence nonvanishing of
\(\vartheta_{a+bc}\) is exactly \(v_{a+bc}\notin q\).
This is the full second clause of
\eqref{eq:ftb-theta-inmm}.
Uniqueness is inherited from projective coordinates and
the open immersion.  Pull the actual line and sections
to \(\mathcal M_0\), which proves the assertion.
\end{proof}

This constructs the source's \(\mathrm{MMPre}\), not yet
its universal abelian scheme.
In particular reducedness and flatness of the coordinate
closure do not by themselves supply any polarized family
on that closure.
The next producer must recover an actual marked theta
object near each relevant prime, prove uniqueness on
overlaps, and glue these objects on a suitable open.
The exact next statements are
\srcid{exists_generic_object},
\srcid{exists_tObj_near}, and
\srcid{iso_unique_of_dense} in the source ledger
below.

\paragraph{Frozen source concordance and quantifiers.}
The definition of
\(\mathcal I,\mathrm{HasTh},I_C,\mathrm{Coset},\mathrm{InMM}\)
is
\srcid{src/Development/AbelianVarietyTheory/Definitions/ModuliSpaces/%
GoodClass.lean}
(\srcid{1f7f9fe959f1a165b471c16211a3b1290fbdeab8}).
Section coordinates and the exact \(\mathrm{MMPre}\)
fields are in
\srcid{FiniteHomogeneousInf.lean}
(\srcid{1a1beedf08bb1a0f2369c22fa441010a9df8d31b})
in the same directory.
The remaining prefix is
\srcid{src/Development/AbelianVarietyTheory/Proofs/ModuliProofs/}.
The actual ideal proof is
\srcid{ThetaPointsVanishingIdealRadicalHomogeneous.lean}
(\srcid{8d8b6ad3c09a9b60f8f95739f59428d949a9f7bc}).
The projective construction and homogeneous chart kernel
are
\srcid{ExistsProjCoordRepresentationHomogeneousIdealFla.lean}
(\srcid{5b0ceb928553b508f4c65e5c8a63dab7f5ceefdc}),
\srcid{ExistsNonzeroCoordinateAffineChart.lean}
(\srcid{dd617586ae8dba48d5326405113f23bdc86eb9b4}),
and
\srcid{DehomogenisationIsRadicalAndTorsionFree.lean}
(\srcid{ab21a5e8d706c8d13251e38044e65adb4db48b96}).
The pullback and stalk nonvanishing comparisons are
\srcid{SectionCoordinatesPullIsoNonzeroLocus.lean}
(\srcid{0799c678994ae1f2ea51199ac6ebd222d6eff3e4}).
The precise assembly is
\srcid{NonemptyThetaCoordinateSchemeData.lean}
(\srcid{1b2990aa07d997e8f510b16185c9a18c28435ebe}).

No analytic separation, generic-object existence, or
polarized extension assertion was needed in this lemma.
Those are separate, substantive next producers:
\srcid{ExistsGenericObjectMinimalPrimeModuliChart.lean}
(\srcid{d82fe58703626f4a4dfe7def4ba0dc2e88b6247e}),
\srcid{ExistsThetaObjectNearGoodPrime.lean}
(\srcid{5694d3c9f8408b96f3cf2535bd140ec9501acea8}),
and
\srcid{ThetaLevelIsoUniqueOfDenseCommonFibres.lean}
(\srcid{1df29797ace199044afe7a78870fd5e45edab58b}).
In the actual fine-family application \(g\) here is \(8h\)
and \(C\) is the eighth-power sharp class.
Its required finite holomorphic lattice-family and
countable-fibre hypotheses, with \(j_0\) chosen before
all arithmetic coefficient rings, remain to be supplied
by the actual sharp-class construction.
